% Part: second-order-logic
% Chapter: metatheory
% Section: compactness

\documentclass[../../../include/open-logic-section]{subfiles}

\begin{document}

\olfileid{sol}{met}{lst}

\olsection{ద్వితీయ-స్థాయి తర్కానికి లొవెన్‌హైమ్--స్కోలెమ్ సిద్ధాంతం విఫలమవుతుంది}

\begin{explain}
అనంత నమూనా గల ప్రతి \tetoken{వాక్యాల}{!!{sentence}s} సమితికీ
\tetoken{ఒక లెక్కించదగిన}{!!a{enumerable}} నమూనా ఉంటుందని (అధోముఖ)
లొవెన్‌హైమ్--స్కోలెమ్ సిద్ధాంతం చెబుతుంది. అది కూడా సంపూర్ణతా సిద్ధాంతపు
పరిణామమే: అవైరుధ్యమైన ఏ \tetoken{వాక్యాల}{!!{sentence}s} సమితికైనా సంపూర్ణత
నిరూపణ ఒక నమూనాను ఉత్పత్తి చేస్తుంది; ఆ నమూనా
\tetoken{లెక్కించదగినది}{!!{enumerable}}. ఒక \tetoken{వాక్యాల}{!!{sentence}s} సమితికి
\tetoken{ఒక అనంతంగా లెక్కించదగిన}{!!a{denumerable}} నమూనా ఉంటే, దానికి
\tetoken{ఒక లెక్కించలేని}{!!a{nonenumerable}} నమూనా కూడా ఉంటుందని హామీ ఇచ్చే
ఊర్ధ్వముఖ లొవెన్‌హైమ్--స్కోలెమ్ సిద్ధాంతం కూడా ఉంది. ఈ రెండు సిద్ధాంతాలూ
ద్వితీయ-స్థాయి తర్కంలో విఫలమవుతాయి.
\end{explain}


\begin{thm}
\ollabel{thm:sol-no-ls} ద్వితీయ-స్థాయి తర్కానికి లొవెన్‌హైమ్--స్కోలెమ్
సిద్ధాంతం విఫలమవుతుంది: అనంత నమూనాలు ఉండి,
\tetoken{లెక్కించదగిన}{!!{enumerable}} నమూనాలు లేని \tetoken{వాక్యాలు}{!!{sentence}s} ఉన్నాయి.
\end{thm}

\begin{proof}
ఈ సూత్రాన్ని గుర్తుచేసుకోండి:
\[
\fn{Count} \ident \lexists[z][\lexists[u][\lforall[X][((X(z) \land
      \lforall[x][(X(x) \lif X(u(x)))]) \lif \lforall[x][X(x)])]]]
\]
ఇది \tetoken{ఒక నిర్మాణం}{!!a{structure}}~$\Struct{M}$లో సత్యమైతేనే
$\Domain{M}$ \tetoken{లెక్కించదగినది}{!!{enumerable}}; అందువల్ల
$\lnot\fn{Count}$,~$\Struct{M}$లో సత్యమైతేనే $\Domain{M}$
\tetoken{లెక్కించలేనిది}{!!{nonenumerable}}. అలాంటి
\tetoken{నిర్మాణాలు}{!!{structure}s} ఉన్నాయి---\tetoken{వ్యక్తి క్షేత్రంగా}{!!{domain}}
ఏదైనా \tetoken{లెక్కించలేని}{!!{nonenumerable}} సమితిని తీసుకోండి; ఉదాహరణకు,
$\Pow{\Nat}$ లేదా~$\Real$. కాబట్టి $\lnot \fn{Count}$కు అనంత నమూనాలు
ఉన్నాయి, కానీ \tetoken{లెక్కించదగిన}{!!{enumerable}} నమూనాలు లేవు.
\end{proof}

\begin{thm}
\tetoken{అనంతంగా లెక్కించదగిన}{!!{denumerable}} నమూనాలు ఉండి,
\tetoken{లెక్కించలేని}{!!{nonenumerable}} నమూనాలు లేని \tetoken{వాక్యాలు}{!!{sentence}s} ఉన్నాయి.
\end{thm}

\begin{proof}
$\fn{Count} \land \fn{Inf}$, $\Nat$లో సత్యం; కానీ $\Domain{M}$
\tetoken{లెక్కించలేనిదిగా}{!!{nonenumerable}} ఉన్న ఏ
\tetoken{నిర్మాణం}{!!{structure}}~$\Struct{M}$లోనూ సత్యం కాదు.
\end{proof}


\end{document}
